\chapter{Calculus I --- Change}
\label{ch:calculus1}

\begin{tcolorbox}[feynbox={Sky}{BBg},title={Before and After}]
\textbf{Before this chapter you need:} functions, domain and range
(\S\ref{sec:functions}--\S\ref{sec:domain}), slope of a straight line
(\S\ref{sec:linear}), quadratics and curvature (\S\ref{sec:quadratics}),
the exponential and the logarithm (\S\ref{sec:compound}--\S\ref{sec:loginverse}).\\
\textbf{After it you will understand:} what a derivative is and how to compute
one, what a second derivative says about the shape of a curve, and what an
integral accumulates. Chapter~\ref{ch:calculus2} lifts all of it to many
variables, which is where machine learning actually lives.
\end{tcolorbox}

Calculus is the mathematics of change. Everything in this thesis that
\emph{learns} does so by measuring how a number responds to a small nudge and
then nudging in the direction that helps. This chapter builds that measurement
for functions of one variable.

%% ─────────────────────────────────────────────────────────────────────────────
\section{Limits, Informally}
\label{sec:limits}

\situation{
A runner covers 100 metres in 10 seconds. Her average speed is 10 m/s, but she
was not travelling at 10 m/s throughout --- she started at zero. To find her speed
at the exact instant $t = 4$ seconds, you might time her over the interval from
4 to 5 seconds. That is still an average. So you shrink the window: 4 to 4.1, then
4 to 4.01, then 4 to 4.001.
}

\problem{
You cannot shrink the window to nothing. Over an interval of exactly zero
duration she covers exactly zero distance, and $0/0$ is not a number. Yet the
sequence of averages over shrinking windows plainly settles down toward a
particular value. We need language for ``the value it settles toward'' that does
not require ever reaching a window of zero width.
}

\workedarith{
Take $f(x) = x^2$ and ask for the slope at $x = 3$.

Average slope from $3$ to $3+h$ is
\[
  \frac{f(3+h) - f(3)}{h} = \frac{(3+h)^2 - 9}{h}
  = \frac{9 + 6h + h^2 - 9}{h} = \frac{6h + h^2}{h} = 6 + h ,
\]
where the cancellation is legal for every $h \neq 0$.

\medskip
\begin{tabular}{@{}lll@{}}
\toprule
$h$ & Average slope $6+h$ \\
\midrule
$1$      & $7$ \\
$0.1$    & $6.1$ \\
$0.01$   & $6.01$ \\
$0.001$  & $6.001$ \\
$-0.001$ & $5.999$ \\
$-0.1$   & $5.9$ \\
\bottomrule
\end{tabular}

\medskip
Approaching from either side, the values close in on $6$. At $h = 0$ the
expression $(6h+h^2)/h$ is $0/0$ and says nothing --- but $6 + h$ at $h = 0$ is
plainly $6$. The limit is $6$.
}

\formaldef{
We write
\[
  \lim_{h \to 0} g(h) = L
\]
to mean: the values $g(h)$ can be made as close to $L$ as desired by taking $h$
sufficiently close to (but not equal to) $0$.

The limit describes the value $g$ \emph{approaches}; whether $g(0)$ exists, and
what it equals if it does, is a separate question. In the example above
$g(h) = (6h+h^2)/h$ is undefined at $h=0$ yet has limit $6$ there.
}

\analogybreak{
Not every function settles. Consider $g(h) = 1/h$ near zero: from the right it
grows without bound, from the left it falls without bound, and there is no value
it approaches. Others approach different values from each side and so have no
limit either. The limit is a claim that must be checked, not a formality --- and
the same care resurfaces in Chapter~\ref{ch:heavytails}, where the sample
kurtosis of a heavy-tailed series has no limit at all, drifting upward without
settling as $n$ grows.
}

\whyhere{
Every derivative in this thesis is a limit of this kind. The autograd machinery
in PyTorch does not compute limits numerically --- it applies the rules of
\S\ref{sec:rules}, which were themselves derived from limits once and for all.
}

%% ─────────────────────────────────────────────────────────────────────────────
\section{The Derivative as Instantaneous Slope}
\label{sec:derivative}

\situation{
Your car's speedometer reads 60 km/h. It is not telling you a distance divided by
a duration over any interval you could name --- it is reporting how fast the
odometer is changing \emph{right now}. The dashboard shows one number, and it is
the limit of the previous section made physical.
}

\problem{
For a straight line, one slope describes the whole thing (\S\ref{sec:linear}).
For anything that bends, the slope differs at every point, so ``the slope'' is not
a number but a new \emph{function} of position. Building that function is the
central construction of calculus.
}

\workedarith{
\textbf{The two-point picture first.} For $f(x) = x^2$, from $x=1$ to $x=4$:
heights $1$ and $16$, change $15$, run $3$, ratio $5$. That is the average slope
over the interval --- the slope of the straight line joining the two points, called
a \emph{secant}.

Shrink the interval toward $x = 1$:

\medskip
\begin{tabular}{@{}lllll@{}}
\toprule
From & To & Rise & Run & Slope \\
\midrule
$1$ & $4$    & $15$      & $3$    & $5$ \\
$1$ & $2$    & $3$       & $1$    & $3$ \\
$1$ & $1.1$  & $0.21$    & $0.1$  & $2.1$ \\
$1$ & $1.01$ & $0.0201$  & $0.01$ & $2.01$ \\
\bottomrule
\end{tabular}

\medskip
Settling on $2$. Repeating the algebra of \S\ref{sec:limits} at a general point:
\[
  \frac{(x+h)^2 - x^2}{h} = \frac{2xh + h^2}{h} = 2x + h
  \;\xrightarrow[h \to 0]{}\; 2x .
\]

So the derivative of $x^2$ is $2x$: at $x=1$ it is $2$, at $x=3$ it is $6$
(matching \S\ref{sec:limits}), at $x=0$ it is $0$ --- the flat bottom of the bowl.

\medskip
\textbf{Sanity check against \S\ref{sec:quadratics}.} There the average slopes
over unit steps were $1, 3, 5, 7$. The derivative at the \emph{midpoints}
$0.5, 1.5, 2.5, 3.5$ is $2x = 1, 3, 5, 7$. Exactly. \checkmark
}

\formaldef{
The \textbf{derivative} of $f$ at $x$ is
\[
  f'(x) = \lim_{h \to 0} \frac{f(x+h) - f(x)}{h} ,
\]
when the limit exists. Equivalent notations: $f'(x)$, $\dfrac{df}{dx}$,
$\dfrac{d}{dx}f(x)$.

$f'(x)$ is the slope of the \textbf{tangent} --- the straight line touching the
curve at $x$ and matching its direction there. A function with a derivative
everywhere on an interval is \textbf{differentiable} there.
}

\analogybreak{
Not every function has a derivative everywhere. The absolute value $|x|$ has slope
$-1$ to the left of zero and $+1$ to the right, and no single tangent at the
corner: the limit from each side disagrees. This is not an exotic pathology. The
ReLU activation used throughout deep learning is exactly this shape, and the
gradient penalty of Chapter~\ref{ch:gans} involves $\bigl|\,\norm{\nabla D} - 1\bigr|$,
which has the same kink. Practical systems patch the corner by picking a
convention; the mathematics genuinely has no answer there.
}

\whyhere{
Training any model in this thesis means computing $\partial(\text{loss})/\partial(\text{weight})$
for every weight and stepping against it. The gradient-clipping threshold of
$1.0$ applied in all three generators is a bound on the size of exactly these
derivatives.
}

%% ─────────────────────────────────────────────────────────────────────────────
\section{The Rules of Differentiation}
\label{sec:rules}

\situation{
Nobody computes limits by hand to differentiate. Once the limit has been
evaluated for a handful of basic functions, and once a few rules for combining
them are established, differentiation becomes bookkeeping --- which is precisely
why a computer can do it automatically.
}

\problem{
The functions we actually differentiate are built by adding, multiplying,
dividing and composing simpler ones. Without combination rules, each new
combination would require returning to the limit definition from scratch.
}

\workedarith{
\textbf{Power rule:} $\dfrac{d}{dx}x^n = n x^{n-1}$.

Check on $x^3$ at $x = 2$. The rule gives $3x^2 = 12$. Numerically, with
$h = 0.001$:
\[
  \frac{(2.001)^3 - 2^3}{0.001} = \frac{8.012006001 - 8}{0.001} = 12.006 ,
\]
converging to $12$ as $h$ shrinks. \checkmark

\medskip
\textbf{Sum rule:} $(f+g)' = f' + g'$. For $f(x) = x^2 + x^3$,
$f'(x) = 2x + 3x^2$. At $x=2$: $4 + 12 = 16$.

\medskip
\textbf{Product rule:} $(fg)' = f'g + fg'$.

Take $f = x^2$, $g = x^3$, so $fg = x^5$ and the power rule predicts $5x^4$.
Via the product rule: $2x \cdot x^3 + x^2 \cdot 3x^2 = 2x^4 + 3x^4 = 5x^4$. \checkmark

Note that $(fg)' \neq f'g'$: here $f'g' = 2x \cdot 3x^2 = 6x^3$, which is wrong.

\medskip
\textbf{Chain rule:} $\bigl(f(g(x))\bigr)' = f'(g(x)) \cdot g'(x)$.

Take $f(u) = u^2$ and $g(x) = 3x + 1$, so $f(g(x)) = (3x+1)^2$. The rule gives
$2(3x+1) \cdot 3 = 6(3x+1) = 18x + 6$. Expanding first:
$(3x+1)^2 = 9x^2 + 6x + 1$, whose derivative is $18x + 6$. \checkmark

\medskip
\textbf{Exponential and logarithm:}
\[
  \frac{d}{dx}e^x = e^x , \qquad \frac{d}{dx}\ln x = \frac{1}{x} \;\; (x>0) .
\]
The first is the property promised in \S\ref{sec:compound}: $e^x$ is its own
derivative, the unique function (up to scaling) that is its own rate of change.
}

\formaldef{
For differentiable $f, g$ and constant $c$:
\[
\begin{aligned}
  (cf)' &= cf' &\qquad (f \pm g)' &= f' \pm g' \\
  (fg)' &= f'g + fg' &\qquad \left(\frac{f}{g}\right)' &= \frac{f'g - fg'}{g^2} \;\; (g \neq 0) \\
  (f \circ g)'(x) &= f'(g(x))\,g'(x) &\qquad \frac{d}{dx}x^n &= nx^{n-1}
\end{aligned}
\]
The chain rule is the one that matters most: it is the mathematical content of
backpropagation.
}

\analogybreak{
The rules make differentiation look purely mechanical, and symbolically it is.
But mechanical is not the same as numerically harmless. Applying the chain rule
through fifty composed layers multiplies fifty derivatives together, and if each
is around $0.5$ the product is about $10^{-15}$ --- the vanishing-gradient problem.
The algebra is exact; the arithmetic that implements it is not.
}

\whyhere{
The chain rule applied to a deep composition is backpropagation, and it is why
Chapter~\ref{ch:calculus2}'s multivariable version is the load-bearing piece of
this whole book. The vanishing-gradient issue above is the direct reason this
project replaced sigmoid activations with $\tanh$ in the TimeGAN latent path:
$\tanh'(0) = 1$ against $\sigma'(0) = 0.25$, a fourfold stronger signal per layer.
}

%% ─────────────────────────────────────────────────────────────────────────────
\section{Second Derivatives and Curvature}
\label{sec:secondderiv}

\situation{
You are in a car. The speedometer says how fast you are going; the push you feel
in your back says how fast \emph{that} is changing. Speed is the first
derivative of position, acceleration the second. They answer different questions,
and you can have one without the other: at constant 100 km/h the speed is large
and the acceleration is zero.
}

\problem{
The first derivative alone cannot tell a valley from a peak. At the bottom of a
bowl and at the top of a hill the slope is zero in both cases. Distinguishing
them requires knowing which way the curve bends.
}

\workedarith{
Take $f(x) = x^2$, so $f'(x) = 2x$ and $f''(x) = 2$.

At $x=0$: $f'(0) = 0$, so it is a turning point. $f''(0) = 2 > 0$, meaning the
slope is increasing as we pass through --- from negative, to zero, to positive.
That is a valley.

Check directly: $f(-1) = 1$, $f(0) = 0$, $f(1) = 1$. Lower in the middle. \checkmark

\medskip
Now $g(x) = -x^2$, so $g'(x) = -2x$ and $g''(x) = -2$.

At $x=0$: $g'(0) = 0$ again, but $g''(0) = -2 < 0$. The slope is decreasing
through the point --- positive, zero, negative. A peak.

Check: $g(-1) = -1$, $g(0) = 0$, $g(1) = -1$. Higher in the middle. \checkmark

\medskip
\textbf{The ambiguous case.} $h(x) = x^3$ has $h'(0) = 0$ and $h''(0) = 0$. It is
neither peak nor valley: $h(-1) = -1$ and $h(1) = 1$, so the function passes
straight through while momentarily flattening. When the second derivative
vanishes, the test is silent.
}

\formaldef{
The \textbf{second derivative} $f''(x)$ is the derivative of $f'$. It measures
\textbf{concavity}:
\begin{itemize}
  \item $f''(x) > 0$: convex here (curving upward, holds water);
  \item $f''(x) < 0$: concave here (curving downward, sheds water);
  \item $f''(x) = 0$: no information from this test.
\end{itemize}
\textbf{Second-derivative test.} If $f'(a) = 0$ and $f''(a) > 0$, then $a$ is a
local minimum; if $f'(a) = 0$ and $f''(a) < 0$, a local maximum.
}

\analogybreak{
``Local'' is doing heavy lifting. The test inspects an infinitesimal neighbourhood
and certifies nothing beyond it: a function may have a local minimum at $a$ and
still take far lower values somewhere else entirely. Every optimisation result in
this thesis is local in exactly this sense --- nothing in the training procedure
certifies a global optimum, and nothing in the reported numbers should be read as
claiming one.
}

\whyhere{
Convexity is what makes an optimisation problem easy, and its absence is what
makes GAN training hard. A convex loss has one basin and any downhill walk finds
it. A GAN's objective is not convex, and moreover the two networks are optimising
against each other, so the surface each is descending shifts as the other learns.
This is the structural reason for the instability Chapter~\ref{ch:gans}
documents, and for the mode collapse the post-generation guard in this project is
built to detect.
}

%% ─────────────────────────────────────────────────────────────────────────────
\section{Integrals as Accumulated Area}
\label{sec:integrals}

\situation{
A tap fills a bucket. The flow rate varies --- fast at first, then slowing. You
know the rate at every instant and want the total volume after ten minutes. If the
rate were constant you would multiply rate by time. It is not, so you cannot.
}

\problem{
Differentiation goes from total to rate. The reverse trip --- rate to total --- is
needed just as often, and it is not simply the same operation run backwards in
any obvious way. It requires its own construction.
}

\workedarith{
\textbf{Constant rate first.} A tap running at 2 litres/minute for 3 minutes
delivers $2 \times 3 = 6$ litres. On a graph of rate against time this is a
rectangle of height 2 and width 3 --- area 6. Total is area under the rate curve.

\medskip
\textbf{Varying rate.} Let the rate be $f(t) = 2t$ litres/minute, from $t=0$ to
$t=3$. Approximate with three rectangles of width 1, using the rate at the right
edge of each:
\[
  f(1) \cdot 1 + f(2) \cdot 1 + f(3) \cdot 1 = 2 + 4 + 6 = 12 .
\]
Too big --- the rate is rising, so the right edge overstates each strip. Using
left edges:
\[
  f(0) + f(1) + f(2) = 0 + 2 + 4 = 6 ,
\]
too small. The truth is between 6 and 12.

\medskip
Refine to six strips of width $0.5$, right edges:
\[
  0.5\bigl(f(0.5) + f(1) + \cdots + f(3)\bigr)
  = 0.5(1 + 2 + 3 + 4 + 5 + 6) = 0.5 \times 21 = 10.5 .
\]
Left edges give $0.5(0+1+2+3+4+5) = 7.5$. The bracket has narrowed from
$[6,12]$ to $[7.5, 10.5]$.

\medskip
\textbf{Exactly.} The region under $f(t) = 2t$ from $0$ to $3$ is a triangle of
base 3 and height 6, so its area is $\tfrac12 \times 3 \times 6 = 9$ litres --- and
$9$ sits inside both brackets. \checkmark
}

\formaldef{
The \textbf{definite integral} of $f$ from $a$ to $b$,
\[
  \int_a^b f(x)\,dx ,
\]
is the limit of such rectangle sums as the widths shrink to zero. It equals the
signed area between the graph and the horizontal axis --- signed because regions
below the axis count negatively.

The elongated $S$ is a stretched ``sum'', and $dx$ records the vanishing width of
each strip.
}

\analogybreak{
Area is the right picture in one dimension and becomes a misleading one quickly.
In the probability chapters an integral over a two-dimensional region is a volume,
and the integrals defining Wasserstein distance are over spaces of \emph{joint
distributions} --- objects with no spatial interpretation at all. What survives
from the picture is the accounting: an integral totals up infinitely many
infinitesimal contributions. That much always holds.
}

\whyhere{
A probability density is a rate --- probability per unit of return --- so
probabilities are integrals of it. The statement that a $5\sigma$ day has
probability $5.73 \times 10^{-7}$ under a Gaussian is the value of an integral over
the two tails, and that single number is what yields the once-per-6{,}922-years
figure that motivates this entire thesis.
}

%% ─────────────────────────────────────────────────────────────────────────────
\section{The Fundamental Theorem}
\label{sec:ftc}

\situation{
Your car has an odometer and a speedometer. The speedometer reading is the rate
of change of the odometer. Equally, the odometer reading is the accumulation of
the speedometer over the journey. Two instruments, one relationship, read in
opposite directions.
}

\problem{
Computing integrals by shrinking rectangles is exhausting and, for most
functions, impossible by hand. If accumulation really is the inverse of
differentiation, then every derivative already known becomes an integral known
for free --- and \S\ref{sec:rules} supplied a whole table of them.
}

\workedarith{
Redo the tap by reversing differentiation. We need a function whose derivative is
$2t$. From the power rule, $\frac{d}{dt}t^2 = 2t$, so $F(t) = t^2$ works.

Then
\[
  \int_0^3 2t\,dt = F(3) - F(0) = 9 - 0 = 9 ,
\]
matching the triangle exactly, with no rectangles at all. \checkmark

\medskip
\textbf{Second example.} $\displaystyle\int_1^4 2x\,dx = F(4) - F(1) = 16 - 1 = 15$.

Geometric check: the region is a trapezium with parallel sides $f(1)=2$ and
$f(4)=8$, width $3$. Area $= \tfrac{1}{2}(2+8) \times 3 = 15$. \checkmark

\medskip
\textbf{Third.} $\displaystyle\int_1^{e} \frac{1}{x}\,dx = \ln(e) - \ln(1) = 1 - 0 = 1$,
since $\frac{d}{dx}\ln x = 1/x$. The area under $1/x$ from $1$ to $e$ is exactly
$1$ --- which is one way of saying what the number $e$ is.
}

\formaldef{
\textbf{Fundamental Theorem of Calculus.} If $F' = f$ on $[a,b]$ and $f$ is
continuous, then
\[
  \int_a^b f(x)\,dx = F(b) - F(a) .
\]
Conversely, $\displaystyle\frac{d}{dx}\int_a^x f(t)\,dt = f(x)$.

$F$ is an \textbf{antiderivative} of $f$. Antiderivatives are not unique --- any
constant may be added, since constants differentiate to zero --- but the constant
cancels in $F(b) - F(a)$, so the definite integral is unaffected.
}

\analogybreak{
The theorem guarantees an antiderivative exists for any continuous $f$; it does
not guarantee you can write one down. The Gaussian density $e^{-x^2/2}$ has no
antiderivative expressible in elementary functions, which is exactly why the
normal CDF $\Phi$ must be tabulated or computed numerically. The
$5.73 \times 10^{-7}$ quoted above is a numerical result, not a closed-form one.
}

\whyhere{
The theorem links densities to probabilities throughout this thesis: the CDF is
the integral of the density, and the density is the derivative of the CDF. Every
quantile --- and Value at Risk is precisely a quantile --- is found by inverting
that relationship.
}
